\subsection{Real-Telemetry Detection and Fault-Domain Diagnosis}

The audited UAV-SEAD corpus contained 1396 mapped logs, of which 1389 were retained as usable after excluding missing, unparsable, or corrupt records. The retained logs yielded 1,892,063 aligned records at 10 Hz and 87 core telemetry channels. Applying the fixed AeroTSBoost-compatible representation produced 218,537 overlapping windows with 1566 temporal-statistical features per window, including 15,694 anomalous windows. Under the fixed per-log chronological split, 152,340, 32,138, and 34,059 windows were assigned to training, validation, and test partitions, respectively; the corresponding anomalous-window counts were 9011, 3081, and 3602. These audit values matched the expected dataset and representation counts before model evaluation.

The five-seed LightGBM detector reproduced the reference performance level on the fixed chronological test partition. The mean AUROC was $0.9512 \pm 0.0006$ and the mean AUPRC was $0.7522 \pm 0.0039$. The independently optimized test-set F1 was $0.6924 \pm 0.0056$. These values were close to the reference values of 0.9516, 0.7516, and 0.6860, respectively, and preserved the reported detector performance trend. The AUPRC result is the primary threshold-independent measure because anomalous windows were a minority of the evaluated windows.

For operational use, the threshold was selected separately for each seed by maximizing validation F1 and then frozen before test evaluation. With this validation-derived threshold, the test F1 was $0.6895 \pm 0.0067$, and the log-aware event F1 was $0.6051 \pm 0.0218$. Thus, the threshold-dependent operating result remained close to the independently optimized test F1, while the event-level measure was lower because it evaluates the temporal detection behavior at the frozen operating point.

On the 3602 anomalous test windows, the anomaly-only Stage 2 LightGBM classifier attained Macro-F1 of $0.8980 \pm 0.0020$ and Balanced Accuracy of $0.8863 \pm 0.0021$. Its mean per-class recall and F1 were 0.9685 and 0.9390 for External Position, 0.9279 and 0.9162 for Global Position, 0.7916 and 0.8503 for Altitude, and 0.8574 and 0.8865 for Mechanical/Electrical, respectively. The corresponding Random Forest baseline achieved Macro-F1 of $0.8724 \pm 0.0037$ and Balanced Accuracy of $0.8382 \pm 0.0052$. The majority and stratified-random references achieved Macro-F1 values of 0.1608 and $0.2358 \pm 0.0090$. These comparisons indicate that the four-class diagnostic task was learned under the specified anomaly-only label protocol; they do not represent complete online five-class performance.

When Stage 1 thresholding and Stage 2 prediction were combined, the complete five-class Cascade achieved Macro-F1 of $0.6902 \pm 0.0071$ and Balanced Accuracy of $0.6755 \pm 0.0152$. Its mean recall/F1 pairs were 0.9629/0.9632 for Normal, 0.7787/0.7088 for External Position, 0.6649/0.6493 for Global Position, 0.4329/0.5186 for Altitude, and 0.5379/0.6112 for Mechanical/Electrical. The lower anomaly-class performance relative to the conditional Stage 2 result reflects the additional Stage 1 gating step. The five-seed Direct Five-Class LightGBM comparator achieved Macro-F1 of $0.6599 \pm 0.0029$ and Balanced Accuracy of $0.6144 \pm 0.0021$ on the same fixed chronological test partition. These fixed-split values describe the primary reproduction protocol; stricter isolation results are reported separately.
